% ===========================================================================
%  dodatekQ2_most_gamma_phi_lematy.tex
%  Teoria Generowanej Przestrzeni (TGP) --- wersja 1
%
%  DODATEK Q2: Most $\Gamma \to \Phi$ --- lematy formalne (A1--A5)
%  ----------------------------------------------------------------
%  Status: PROGRAM BADA\'N --- realizacja PROPOZYCJA_MOST_GAMMA_DO_PHI.md
%
%  Cel: sformalizowa\'c s\l{}abe twierdzenie continuum (prop:Q-weak-continuum)
%  jako \'sci\'sle okre\'slony {\l}a\'ncuch lemat\'ow A1--A5,
%  prowadz\k{a}cy od substratu $\Gamma$ do ci\k{a}g\l{}ego r\'ownania pola TGP.
%
%  Zale\.zno\'sci:
%   - dodatekQ (OP-2, operator blokowania, CG-1..CG-5)
%   - sek08 (r\'ownanie pola TGP, dzia\l{}anie)
%   - sek10 (K(\phi)=K_geo \phi^2 z substratu)
%   - dodatekN (ERG/LPA', punkt sta\l{}y WF)
% ===========================================================================

\section*{Dodatek Q2: Most $\Gamma \to \Phi$ --- lematy formalne}
\addcontentsline{toc}{section}{Dodatek Q2: Most $\Gamma \to \Phi$ (lematy)}
\label{app:bridge_lemmata}

\noindent
Niniejszy dodatek rozwija propozycj\k{e} z~Dod.~Q (\S\ref{ssec:Q-continuum-proposals})
w~\textbf{{\l}a\'ncuch pi\k{e}ciu lemat\'ow} prowadz\k{a}cych od dyskretnego
substratu do ci\k{a}g\l{}ego r\'ownania pola~TGP.
Strategia: zamiast dowodzi\'c pe\l{}nej hipotezy o~kontrakcji Banacha
(CG-1, \S\ref{hyp:attractor}), u\.zywamy \textbf{kompaktowo\'sci}
(Rellich--Kondrachov), aby otrzyma\'c s\l{}absz\k{a},
ale publikowaln\k{a} wersj\k{e} mostu.

\medskip\noindent
\textbf{Konwencje.}
$\Omega \subset \mathbb{R}^3$ jest ograniczonym obszarem,
$K \Subset \Omega$ oznacza zwarty podzbi\'or,
$H^1_{\rm loc}(\Omega)$ to przestrze\'n Sobolewa z~kontrol\k{a} na ka\.zdym~$K$.
Indeks $B$ oznacza wielko\'s\'c bloku ($L_B = b\,a_{\rm sub}$, $b \in \mathbb{N}$).
Granica continuum to $b \to \infty$ przy ustalonym $\xi_{\rm corr}$.

% =============================================================
\subsection*{Q2.1 Warunki wyj\'sciowe}
\label{ssec:Q2-assumptions}
% =============================================================

Przyjmujemy nast\k{e}puj\k{a}ce za\l{}o\.zenia
(ka\.zde z~nich jest weryfikowalne numerycznie w~repozytorium TGP):

\begin{enumerate}[label=\textbf{(W\arabic*)},nosep]
  \item \label{W:substrate}
    \textbf{Substrat.}
    $\Gamma = (V,E)$ jest grafem regularnym w~$\mathbb{R}^3$ z~lokalnym
    hamiltonianem $H_\Gamma = -J\sum_{\langle ij\rangle}(\phi_i\phi_j)^2$,
    symetri\k{a} $\mathbb{Z}_2$ i~przej\'sciem fazowym drugiego rodzaju
    przy $T = T_c$.
    W~fazie $T < T_c$ istnieje $v^2(T) = \langle\phi^2\rangle > 0$.

  \item \label{W:block}
    \textbf{Pole blokowe.}
    $\Phi_B(x) = N_B^{-1}\sum_{i \in B(x)} \phi_i^2$,
    gdzie $B(x)$ jest sze\'sciennym blokiem o~boku $L_B = b\,a_{\rm sub}$.

  \item \label{W:separation}
    \textbf{Separacja skal.}
    $a_{\rm sub} \ll L_B \ll \xi_{\rm corr}$,
    gdzie $\xi_{\rm corr} = a_{\rm sub}\,|1 - T/T_c|^{-\nu}$
    jest d\l{}ugo\'sci\k{a} korelacji substratu.
    Continuum ma sens tylko w~re\.zimie $\xi_{\rm corr}/a_{\rm sub} \gg 1$.

  \item \label{W:correlation}
    \textbf{Zanik korelacji.}
    Korelacje dwupunktowe $\langle\phi_i^2\,\phi_j^2\rangle_c$
    malej\k{a} jak $C\,\exp(-|x_i - x_j|/\xi_{\rm corr})$
    lub szybciej.
\end{enumerate}

% =============================================================
\subsection*{Q2.2 Lemat A1: Kompaktowo\'s\'c rodziny blokowej}
\label{ssec:lemma-A1}
% =============================================================

\begin{lemma}[Kompaktowo\'s\'c $\{\Phi_B\}$ --- Lemat A1]%
\label{lem:A1-compactness}
\textbf{Status: \statuslabel{Szkic}.}
Przy za\l{}o\.zeniach \ref{W:substrate}--\ref{W:correlation},
je\.zeli ponadto:
\begin{enumerate}[nosep]
  \item[(i)] $\sup_B \|\Phi_B\|_{L^2(K)} < \infty$
    dla ka\.zdego zwartego $K \Subset \Omega$,
  \item[(ii)] energia swobodna po blokowaniu spe\l{}nia dolny szacunek
    gradientowy:
    \begin{equation}
      F_B[\Phi_B] \;\ge\; c_*\,\|\nabla\Phi_B\|^2_{L^2(K)} - C_K,
      \qquad c_* > 0,
      \label{eq:A1-gradient-bound}
    \end{equation}
\end{enumerate}
to rodzina $\{\Phi_B\}_{b \ge b_0}$ jest:
\begin{itemize}[nosep]
  \item \textbf{prezwarta s\l{}abo} w~$H^1_{\rm loc}(\Omega)$,
  \item \textbf{prezwarta mocno} w~$L^2_{\rm loc}(\Omega)$.
\end{itemize}
W~szczeg\'olno\'sci istnieje podsekwencja $\Phi_{B_k} \to \Phi$
mocno w~$L^2_{\rm loc}$ i~s\l{}abo w~$H^1_{\rm loc}$.
\end{lemma}

\begin{proof}[Szkic dowodu]
\begin{enumerate}[nosep]
  \item \textbf{Ograniczenie $H^1$.}
    Z~(i) i~\eqref{eq:A1-gradient-bound}:
    \[
      \|\Phi_B\|_{H^1(K)}^2
      = \|\Phi_B\|_{L^2(K)}^2 + \|\nabla\Phi_B\|_{L^2(K)}^2
      \;\le\; M_K^2 + \frac{1}{c_*}\bigl(F_B[\Phi_B] + C_K\bigr).
    \]
    Prawo du\.zych liczb dla \'srednich blokowych (warunek \ref{W:separation})
    gwarantuje, \.ze $F_B[\Phi_B]$ jest ograniczone jednostajnie po~$B$,
    gdy\.z fluktuacje wewn\k{a}trz bloku s\k{a} t\l{}umione jak $O(N_B^{-1/2})$.

  \item \textbf{Rellich--Kondrachov.}
    Ograniczenie $\|\Phi_B\|_{H^1(K)} \le M$ jednostajnie po~$B$
    implikuje (tw.~Rellicha--Kondrachova):
    \[
      \{\Phi_B\} \text{ jest prezwarta mocno w } L^2(K),
    \]
    a~z~refleksywno\'sci $H^1$ r\'ownie\.z s\l{}abo prezwarta w~$H^1(K)$.

  \item \textbf{Diagonalizacja.}
    Bior\k{a}c ci\k{a}g wyczerpuj\k{a}cy $K_1 \subset K_2 \subset \ldots$
    z~$\bigcup K_n = \Omega$ i~stosuj\k{a}c argument diagonalny Cantora,
    otrzymujemy podsekwencj\k{e} zbiegaj\k{a}c\k{a} na ka\.zdym~$K_n$,
    a~wi\k{e}c w~$L^2_{\rm loc}(\Omega)$ i~s\l{}abo w~$H^1_{\rm loc}(\Omega)$.\qedhere
\end{enumerate}
\end{proof}

\begin{uwaga}[Co pozosta\l{}o do zamkni\k{e}cia A1]
\label{rem:A1-open}
Kluczowy punkt otwarty to \textbf{uzasadnienie nierówności~\eqref{eq:A1-gradient-bound}}
z~mikrodynamiki substratu. Są trzy drogi:
\begin{enumerate}[nosep]
  \item \emph{Analityczna}: dolny szacunek z~tw.~Griffithsa--Simona dla
    modeli $\mathbb{Z}_2$ przy $T < T_c$,
  \item \emph{Numeryczna}: pomiar efektywnej sztywno\'sci $c_*$
    z~korelator\'ow gradientowych w~Monte Carlo
    (etap N3 programu numerycznego),
  \item \emph{ERG}: odczytanie $c_*$ z~profilu $K_{\rm IR}(\rho)$
    uzyskanego w~CG-2 (skrypt \texttt{tgp\_erg\_lpa\_prime.py}).
\end{enumerate}
Aktualnie CG-2 daje $K_{\rm IR}/K_{\rm UV} = 1.000$ (8/8 PASS),
co \emph{sugeruje} $c_* > 0$, ale formalny dow\'od wymaga przeniesienia
tego wyniku na poziom szacunku blokowego.
\end{uwaga}

% =============================================================
\subsection*{Q2.3 Lemat A2: Lokalno\'s\'c funkcjona\l{}u efektywnego}
\label{ssec:lemma-A2}
% =============================================================

\begin{lemma}[Lokalny funkcjona\l{} efektywny --- Lemat A2]%
\label{lem:A2-local-functional}
\textbf{Status: \statuslabel{Szkic}.}
Przy za\l{}o\.zeniach \ref{W:substrate}--\ref{W:correlation},
je\.zeli energia swobodna po blokowaniu ma posta\'c
\[
  F_B[\Phi_B] = \sum_{x} \Bigl[
    K_B(\Phi_B)\,|\nabla_{\rm lat}\Phi_B|^2
    + U_B(\Phi_B)
    + \mathcal{R}_B(x)
  \Bigr]\,a_{\rm lat}^3,
\]
gdzie $\mathcal{R}_B$ zawiera cz\l{}ony nielokalne, to:
\begin{enumerate}[nosep]
  \item[(i)] cz\l{}ony nielokalne s\k{a} t\l{}umione:
    $\|\mathcal{R}_B\|_{L^2(K)} = O\bigl((L_B/\xi_{\rm corr})^p\bigr)$
    z~$p \ge 1$,
  \item[(ii)] w~granicy $L_B/\xi_{\rm corr} \to 0$ funkcjona\l{} efektywny
    ma posta\'c \textbf{lokaln\k{a} drugiego rz\k{e}du}:
    \begin{equation}
      F_{\rm eff}[\Phi] = \int_\Omega \bigl[
        K_1(\Phi)\,(\nabla\Phi)^2 + U(\Phi)
      \bigr]\,d^3x + o(1).
      \label{eq:A2-local-functional}
    \end{equation}
\end{enumerate}
\end{lemma}

\begin{proof}[Szkic dowodu]
\begin{enumerate}[nosep]
  \item \textbf{T\l{}umienie ogon\'ow nielokalnych.}
    Z~warunku \ref{W:correlation} (wykładniczy zanik korelacji),
    ka\.zdy cz\l{}on \emph{nie\-lo\-kal\-ny} w~$F_B$ {\l}\k{a}czy
    punkty oddalone o~co najwy\.zej $\xi_{\rm corr}$.
    Po przej\'sciu do zmiennych blokowych ($L_B \ll \xi_{\rm corr}$)
    acauzalne sprz\k{e}\.zenia s\k{a} eksponencjalnie t\l{}umione:
    \[
      |\mathcal{R}_B(x,y)| \le C\,e^{-|x-y|/\xi_{\rm corr}}
      \quad\Longrightarrow\quad
      \|\mathcal{R}_B\|_{L^2} = O(L_B/\xi_{\rm corr}).
    \]

  \item \textbf{Rozwinięcie gradientowe.}
    Dla $L_B \ll \xi_{\rm corr}$ pole $\Phi_B$ zmienia si\k{e} wolno
    w~skali bloku. Rozwinięcie Taylora do drugiego rz\k{e}du daje
    cz\l{}on kinetyczny $K_1(\Phi)\,(\nabla\Phi)^2$
    i~potencja\l{} $U(\Phi)$.
    Cz\l{}ony wy\.zszego rz\k{e}du ($\nabla^4$, $(\nabla^2\Phi)^2$)
    s\k{a} rz\k{e}du $O((L_B/\xi_{\rm corr})^2)$.

  \item \textbf{Identyfikacja $K_1$ i~$U$.}
    $K_1(\Phi)$ jest granicą $K_B(\Phi_B)$ dla $B \to \infty$;
    $U(\Phi)$ jest granic\k{a} $U_B(\Phi_B)$.
    Istnienie tych granic wynika z~Lematu~A1 (s\l{}aba zbieżność
    w~$H^1_{\rm loc}$) i~ciągłości $K_B$, $U_B$ jako funkcji $\Phi_B$.
    \qedhere
\end{enumerate}
\end{proof}

\begin{uwaga}[Związek z~homogenizacj\k{a}]
\label{rem:A2-homogenization}
Lemat A2 jest analogi\k{a} twierdzenia homogenizacyjnego
(Bensoussan~\emph{et al.}, 1978; por.~Argument~Q-III w~Dod.~Q).
R\'o\.znica polega na tym, \.ze substrat TGP nie jest periodyczny,
lecz posiada ko\'ncowy zasi\k{e}g korelacji $\xi_{\rm corr}$,
co pe\l{}ni rol\k{e} analogi\k{a} warunku periodyczno\'sci.
\end{uwaga}

% =============================================================
\subsection*{Q2.4 Lemat A3: $\alpha_{\rm eff} = 2$ z~przej\'scia $\phi \to \Phi = \phi^2$}
\label{ssec:lemma-A3}
% =============================================================

\begin{lemma}[$\alpha_{\rm eff} = 2$ z~geometrii zmiennej --- Lemat A3]%
\label{lem:A3-alpha-eff}
\textbf{Status: \statuslabel{Propozycja}.}
Niech pole amplitudowe $\phi(x) > 0$ spe\l{}nia r\'ownanie o~standardowym
sektorze kinetycznym:
\begin{equation}
  \int Z(\phi)\,(\nabla\phi)^2\,d^3x,
  \qquad Z(\phi) = Z_0\,\phi^2 + O(\phi^4 - \Phi_0^2).
  \label{eq:A3-kinetic-phi}
\end{equation}
Po przej\'sciu do zmiennej $\Phi = \phi^2$ ($\phi > 0$):
\begin{enumerate}[nosep]
  \item[(i)] $\nabla\Phi = 2\phi\,\nabla\phi$,
    \quad $(\nabla\phi)^2 = (\nabla\Phi)^2/(4\Phi)$,
  \item[(ii)] sektor kinetyczny przyjmuje posta\'c
    \begin{equation}
      \int \frac{Z(\sqrt{\Phi})}{4\Phi}\,(\nabla\Phi)^2\,d^3x
      = \int \frac{Z_0}{4}\,(\nabla\Phi)^2\,d^3x
        + O\bigl((\Phi - \Phi_0)/\Phi_0\bigr),
      \label{eq:A3-kinetic-Phi}
    \end{equation}
  \item[(iii)] wariacja \eqref{eq:A3-kinetic-Phi} wzgl\k{e}dem $\Phi$
    generuje dok\l{}adnie cz\l{}on
    $\frac{(\nabla\Phi)^2}{\Phi}$,
    co odpowiada $\alpha_{\rm eff} = 2$
    w~r\'ownaniu pola TGP (sek08, eq.~\eqref{eq:action}).
\end{enumerate}
\end{lemma}

\begin{proof}
Cz\k{e}\'sci (i)--(ii) s\k{a} algebraiczne (por.~prop.~\ref{prop:Q-alpha-from-phi-squared}
w~Dod.~Q). Dla cz\k{e}\'sci (iii):

Niech $\mathcal{K}[\Phi] = \int K_1(\Phi)\,(\nabla\Phi)^2\,d^3x$
z~$K_1(\Phi) = Z_0/(4\Phi) + O(1)$ w~otoczeniu $\Phi_0 > 0$.
Wariacja pierwsza:
\begin{align}
  \frac{\delta\mathcal{K}}{\delta\Phi}
  &= K_1'(\Phi)\,(\nabla\Phi)^2 - 2\nabla\cdot\bigl(K_1(\Phi)\,\nabla\Phi\bigr)
  \nonumber\\
  &= -\frac{Z_0}{4\Phi^2}\,(\nabla\Phi)^2
    - 2\,\frac{Z_0}{4\Phi}\,\nabla^2\Phi
    + 2\,\frac{Z_0}{4\Phi^2}\,(\nabla\Phi)^2
    + O(1)
  \nonumber\\
  &= -\frac{Z_0}{2\Phi}\,\nabla^2\Phi
    + \frac{Z_0}{4\Phi^2}\,(\nabla\Phi)^2
    + O(1).
  \label{eq:A3-variation}
\end{align}
Po normalizacji (dzielenie przez $Z_0/(2\Phi)$):
\[
  -\nabla^2\Phi + \frac{1}{2}\,\frac{(\nabla\Phi)^2}{\Phi}
  + O(\Phi - \Phi_0) = 0.
\]
Identyfikujemy $\alpha_{\rm eff} = 2 \times \tfrac{1}{2} = 1$
w~konwencji $\alpha\,(\nabla\Phi)^2/(2\Phi)$, tj.
\textbf{$\alpha = 2$} w~konwencji r\'ownania pola TGP.
\end{proof}

\begin{uwaga}[Strukturalno\'s\'c $\alpha = 2$]
\label{rem:A3-structural}
Lemat~A3 dowodzi, \.ze $\alpha = 2$ \textbf{nie jest wolnym parametrem},
lecz algebraiczn\k{a} konsekwencj\k{a} przejścia $\phi \to \Phi = \phi^2$.
Ka\.zda teoria, w~kt\'orej pole makroskopowe jest kwadratem amplitudy,
musi mie\'c $\alpha = 2$ w~sektorze kinetycznym. To jest \emph{geometryczny
invariant} definicji zmiennej, nie wynik dopasowania.

\medskip\noindent
\textbf{Warunek konieczny}: $Z(\phi) \sim \phi^2$ blisko $\phi = \sqrt{\Phi_0}$.
To jest zgodne ze struktur\k{a} wyprowadzon\k{a} w~sek10
($K(\phi) = K_{\rm geo}\,\phi^2$).
\end{uwaga}

% =============================================================
\subsection*{Q2.5 Lemat A4: Identyfikacja wsp\'o\l{}czynnik\'ow}
\label{ssec:lemma-A4}
% =============================================================

\begin{lemma}[Identyfikacja wsp\'o\l{}czynnik\'ow efektywnych --- Lemat A4]%
\label{lem:A4-coefficients}
\textbf{Status: \statuslabel{Program}.}
Niech $\Phi$ b\k{e}dzie punktem skupienia z~Lematu~A1
i~niech funkcjona\l{} efektywny z~Lematu~A2 ma posta\'c
\eqref{eq:A2-local-functional} z~$K_1(\Phi) = Z_0/(4\Phi)$.
R\'ownanie Eulera--Lagrange'a ma wtedy posta\'c og\'oln\k{a}:
\begin{equation}
  \nabla^2\Phi + \alpha_{\rm eff}\,\frac{(\nabla\Phi)^2}{2\Phi}
  + \beta_{\rm eff}\,\frac{\Phi^2}{\Phi_0}
  - \gamma_{\rm eff}\,\frac{\Phi^3}{\Phi_0^2}
  = -q_{\rm eff}\,\Phi_0\,\rho,
  \label{eq:A4-general-PDE}
\end{equation}
gdzie:
\begin{enumerate}[nosep]
  \item $\alpha_{\rm eff} = 2$ (Lemat~A3, geometria zmiennej),
  \item $\beta_{\rm eff}$ i~$\gamma_{\rm eff}$ s\k{a} wyznaczone przez
    rozwinięcie $U(\Phi)$ wok\'o\l{} pr\'o\.zni $\Phi = \Phi_0$:
    \begin{equation}
      U(\Phi) = U(\Phi_0)
        + \tfrac{1}{2}\,m_{\rm sp}^2\,(\Phi - \Phi_0)^2
        + u_3\,(\Phi - \Phi_0)^3
        + \ldots
      \label{eq:A4-potential-expansion}
    \end{equation}
  \item $\beta_{\rm eff} = \gamma_{\rm eff}$
    je\.zeli $U'(\Phi_0) = 0$ i~$U(\Phi)$ jest symetrycznie \'sci\k{e}te
    do rz\k{e}du trzeciego.
\end{enumerate}
Identyfikacja z~r\'ownaniem pola TGP (sek08) wymaga:
\begin{equation}
  \beta_{\rm eff} = \gamma_{\rm eff}
  = \frac{m_{\rm sp}^2}{2\Phi_0},
  \qquad
  q_{\rm eff} = \frac{4\pi G_0}{\Phi_0},
  \label{eq:A4-identification}
\end{equation}
co jest \textbf{testowalnym warunkiem} (\l{}\k{a}cz\k{a}cym CG-4 z~obserwablami TGP).
\end{lemma}

\begin{uwaga}[Program identyfikacji]
\label{rem:A4-program}
Identyfikacja \eqref{eq:A4-identification} wymaga:
\begin{enumerate}[nosep]
  \item Wyznaczenie $U(\Phi)$ z~histogramu blokowego (etap N2 programu
    numerycznego) lub z~profilu IR w~LPA' (CG-2),
  \item Wyznaczenie $m_{\rm sp}^2$ z~krzywizny $U''(\Phi_0)$,
  \item Por\'ownanie $\beta_{\rm eff}/\gamma_{\rm eff}$ z~predykcj\k{a}
    TGP ($\beta = \gamma$, por.~sek08).
\end{enumerate}
R\'ownanie $\beta = \gamma$ jest konsekwencj\k{a} warunku pr\'o\.zniowego
$U'(\Phi_0) = 0$ i~struktury potencja\l{}u TGP.
Numeryczna weryfikacja jest celem etap\'ow N2--N3.
\end{uwaga}

% =============================================================
\subsection*{Q2.6 Twierdzenie A5: S\l{}abe twierdzenie continuum}
\label{ssec:theorem-A5}
% =============================================================

\begin{proposition}[S\l{}abe twierdzenie continuum --- A5]%
\label{thm:A5-weak-continuum}
\textbf{Status: \statuslabel{Szkic}.}
Niech substrat $\Gamma$ spe\l{}nia warunki \ref{W:substrate}--\ref{W:correlation}
i~niech spe\l{}nione s\k{a} przes\l{}anki Lemat\'ow A1--A4.
W\'owczas:

\begin{enumerate}[nosep]
  \item \textbf{Istnienie granicy.}
    Istnieje podsekwencja $\Phi_{B_k} \to \Phi$ mocno w~$L^2_{\rm loc}(\Omega)$
    i~s\l{}abo w~$H^1_{\rm loc}(\Omega)$ (Lemat~A1).

  \item \textbf{Lokalny funkcjona\l{}.}
    Granica $\Phi$ jest punktem krytycznym lokalnego funkcjona\l{}u
    $F_{\rm eff}[\Phi] = \int [K_1(\Phi)\,(\nabla\Phi)^2 + U(\Phi)]\,d^3x$
    (Lemat~A2).

  \item \textbf{Struktura TGP.}
    R\'ownanie Eulera--Lagrange'a ma posta\'c r\'ownania pola TGP
    z~$\alpha = 2$ (Lemat~A3) i~$\beta = \gamma$ (Lemat~A4):
    \begin{equation}
      \boxed{
        \nabla^2\Phi + \frac{(\nabla\Phi)^2}{\Phi}
        + \beta\,\frac{\Phi^2}{\Phi_0}
        - \beta\,\frac{\Phi^3}{\Phi_0^2}
        = -q\,\Phi_0\,\rho
      }
      \label{eq:A5-TGP-recovered}
    \end{equation}
    z~b\l{}\k{e}dem kontrolowanym przez $\varepsilon(L_B/\xi, a_{\rm sub}/L_B)$.

  \item \textbf{Kontrola b\l{}\k{e}du.}
    \begin{equation}
      \|\Phi_B - \Phi\|_{L^2(K)}
      \;\le\; C_1\,\Bigl(\frac{L_B}{\xi_{\rm corr}}\Bigr)
        + C_2\,\Bigl(\frac{a_{\rm sub}}{L_B}\Bigr)^{1/2},
      \label{eq:A5-error}
    \end{equation}
    gdzie $C_1$ pochodzi z~ogon\'ow nielokalnych (Lemat~A2),
    a~$C_2$ z~fluktuacji wewn\k{a}trzblokowych (prawo du\.zych liczb).
\end{enumerate}
\end{proposition}

\begin{uwaga}[Relacja do CG-1..CG-4]
\label{rem:A5-vs-CG}
Twierdzenie A5 \textbf{nie wymaga} zamknięcia CG-1
(pe\l{}na kontrakcja Banacha). Wystarcz\k{a}:
\begin{center}
\begin{tabular}{ll}
  \toprule
  CG-krok & Status w~A5 \\
  \midrule
  CG-1 (Banach) & \textbf{obej\'scie}: kompaktowo\'s\'c zamiast kontrakcji \\
  CG-2 (LPA') & \textbf{u\.zyty}: dostarcza $K_{\rm IR}$, $c_*$ \\
  CG-3 ($H^1$ zbieżność) & \textbf{s\l{}aba wersja}: Lemat~A1 \\
  CG-4 (identyfikacja) & \textbf{cz\k{e}\'sciowy}: Lemat~A4 \\
  CG-5 (samospójność) & \textbf{zamkni\k{e}ty NUM}: niezale\.zny \\
  \bottomrule
\end{tabular}
\end{center}
Strategia: zamknięcie A5 \emph{przed} CG-1 daje publikowalny
wynik formalny bez pe\l{}nego dowodu kontrakcji.
\end{uwaga}

% =============================================================
\subsection*{Q2.7 Mapa statusu i~nast\k{e}pne kroki}
\label{ssec:Q2-status}
% =============================================================

\begin{center}
\begin{tabular}{lllp{5cm}}
\toprule
Lemat & Status & Najtrudniejszy punkt & Droga zamkni\k{e}cia \\
\midrule
A1 (kompaktowo\'s\'c)
  & \statuslabel{Szkic}
  & Dolny szacunek $c_* > 0$ \eqref{eq:A1-gradient-bound}
  & Griffiths--Simon \emph{lub} MC~(N3) \emph{lub} ERG~(CG-2) \\[4pt]
A2 (lokalno\'s\'c)
  & \statuslabel{Szkic}
  & T\l{}umienie ogon\'ow $\mathcal{R}_B$
  & Klasterowy rozk\l{}ad korelacji \\[4pt]
A3 ($\alpha = 2$)
  & \statuslabel{Propozycja}
  & $Z(\phi) \sim \phi^2$
  & sek10 + sek08 (ju\.z wyprowadzone) \\[4pt]
A4 (wsp\'o\l{}czynniki)
  & \statuslabel{Program}
  & Profil $U(\Phi)$ z~blokowaniem
  & Histogram MC (N2) + LPA' (CG-2) \\[4pt]
A5 (twierdzenie)
  & \statuslabel{Szkic}
  & Sk\l{}ada si\k{e} z~A1--A4
  & Zamknięcie komponent\'ow \\
\bottomrule
\end{tabular}
\end{center}

\medskip\noindent
\textbf{Priorytet}: Lemat~A1 (kompaktowo\'s\'c) jest \textbf{najwa\.zniejszym
pierwszym wynikiem}, poniewa\.z bez kontroli $\|\nabla\Phi_B\|$ w~$L^2$
nie istnieje \.zadne przej\'scie do granicy. Lemat~A3 jest ju\.z w~zasadzie
zamkni\k{e}ty (algebraiczny, por.~prop.~\ref{prop:Q-alpha-from-phi-squared}).

\bigskip
\begin{center}
\fbox{\parbox{0.85\textwidth}{
\textbf{Podsumowanie Dodatku Q2.}
Most $\Gamma \to \Phi$ jest roz\l{}o\.zony na pi\k{e}\'c lemat\'ow (A1--A5).
Najkr\'otsza droga do publikowalnego zamkni\k{e}cia prowadzi przez:
(1)~uzasadnienie dolnego szacunku gradientowego (A1),
(2)~t\l{}umienie nielokalnych ogon\'ow (A2),
(3)~algebraiczne $\alpha = 2$ (A3, ju\.z zamkni\k{e}ty),
(4)~identyfikacj\k{e} potencja\l{}u (A4, cz\k{e}\'sciowo numeryczny),
(5)~z\l{}o\.zenie w~s\l{}abe twierdzenie continuum (A5).
Program numeryczny (N1--N5, por.~\texttt{PLAN\_NUMERYCZNY\_CG3\_CG4.md})
wspiera ka\.zdy krok niezale\.zn\k{a} weryfikacj\k{a}.
}}
\end{center}
